(1) 在 \(x=0\) 处振幅 \(Y=A_i(1+r)\)；入射弦横向牵引小块的力是 \(-T\partial_x y(0^-)\)，取定 \(+y\) 的方向可从弦受边界牵引的反作用确定。为使被动端的阻尼吸收率非负，采用小块方程
\[
M\ddot Y+KY=-T\,y_x(0^-,t)
\]
时，以 \(e^{-i\omega t}\) 表示，\( (K-M\omega^2)(1+r)=-iT k(1-r)\)。记 \(Z=Tk/\omega=\sqrt{T\mu}\)、\(D=K-M\omega^2\)，联立
\[
r=\frac{i\omega Z+D}{i\omega Z-D}.
\]
由于分子和分母模相等，故 \(|r|=1\)。注意不同的应力正方向会将反射相位做等价的符号换元，但反射功率必须保持不变。

(2) 无耗散时，入射、反射波平均能流的比值是 \(|r|^2=1\)。反射相位随 \(D/(\omega Z)\) 连续变化，代表端点的弹性储能与惯性储能交替占优。

(3) 当 \(D=0\)，即 \(\omega=\sqrt{K/M}\)，\(r=1\)，端点位移为反节点（近似自由端相位）；当 \(\omega\to0^+\) 且 \(K>0\) 固定，\(r\to-1\)，端点近似固定。质量、刚度极限须与频率极限区分。

(4) 若加入 \(-b\dot Y\)，小块方程改为 \(M\ddot Y+b\dot Y+KY=-T y_x\)，于是 \(D\) 被替换为 \(D-i b\omega\)，
\[
r_b=\frac{D+i\omega(Z-b)}{i\omega(Z+b)-D}.
\]
因此
\[
\mathcal A=1-|r_b|^2=\frac{4bZ\omega^2}{D^2+\omega^2(Z+b)^2}\ge0.
\]
当 \(b\to0^+\) 时 \(\mathcal A\to0\)；当 \(D=0\) 且 \(b=Z\) 时阻抗匹配、无反射。阻尼项消耗的平均功率与散射能流差吻合。
